\clearpage
\Needspace{12\baselineskip}
\section{Experiment tracking and preliminary checks}
\label{app:experiment_details}

\textbf{Logged coordinates.}
The trainer records configuration, dataset sizes, minibatch training loss,
validation loss, LR, elapsed seconds, throughput, and checkpoint events.
W\&B uses \texttt{train/step} by default; token and elapsed-time coordinates
allow alternative axes. The logging call attaches all three coordinates:

\par\noindent
\begin{minipage}{\linewidth}
\begin{lstlisting}[caption={Step, token, and elapsed-time coordinates for experiment metrics.}]
if metrics and run is not None:
    metrics.update(
        {
            "train/step": completed_steps,
            "train/tokens_seen": completed_steps * tokens_per_step,
            "time/elapsed_seconds": perf_counter() - started_at,
        }
    )
    run.log(metrics)
\end{lstlisting}
\end{minipage}\par

\textbf{Timing and recovery.}
MPS/CUDA is synchronized before timing samples. Throughput intervals include
intervening validation and saving; they are not pure kernel timings. An
initial short LR run experienced a nine-minute system-sleep gap; subsequent
Mac runs used \texttt{caffeinate -i}. Checkpoints retain model, optimizer,
architecture, update count, and RNG states. Resuming starts a new W\&B run
and elapsed-time clock. Interrupted probes do not receive a final checkpoint.

\textbf{Preliminary ledger.}
The following checks preceded the completed low-resource LR, batch, and
architecture comparisons. Short screens establish candidate settings or
runtime correctness, not completed-budget convergence.

\begin{table}[htbp]
    \centering
    \small
    \begin{tabular}{@{}>{\raggedright\arraybackslash}p{0.20\textwidth}
        >{\raggedright\arraybackslash}p{0.34\textwidth}
        >{\raggedright\arraybackslash}p{0.39\textwidth}@{}}
        \toprule
        Stage & Configurations tried & Observation / decision \\
        \midrule
        Throughput checks & Batch sizes 8, 16, 32; steps 51--100 measured
        & Approximately 7,880, 8,237, 8,218 tokens/s. Selected batch 16;
          these short checks do not establish convergence quality. \\
        \addlinespace
        Short LR screening & Peaks $3\times10^{-4}$, $10^{-3}$,
          $3\times10^{-3}$; 1,000 steps each
        & Final validation losses 2.6617, 2.3287, 2.5064.
          Selected $10^{-3}$ for 40.96M-token runs. \\
        \addlinespace
        Architecture smoke checks & Baseline, no RMSNorm, post-norm,
          NoPE, SiLU; 30 MPS steps each
        & All five passed training, validation, checkpoint checks,
          and reconstruction/forward checks. These are runtime checks,
          not formal ablation results. \\
        \addlinespace
        Full LR comparison & Peaks $5\times10^{-4}$, $10^{-3}$,
          $2\times10^{-3}$; 10,000 steps each
        & Final validation losses 1.6402, 1.5966, 1.6092.
          Retained $10^{-3}$ as the best tested setting at 40.96M tokens. \\
        \addlinespace
        Early stability probes & Peaks 0.01, 0.03, 0.1, 0.3, 0.5;
          about 1,000 steps with a 10,000-step schedule
        & Increasingly poor learning; repeated loss explosions at
          0.3 and 0.5, with finite values and intermittent recovery.
          Stopped probes manually. \\
        \bottomrule
    \end{tabular}
    \caption{Completed checks and experiments. All LR experiments used seed 42;
    minimum LR was one tenth of peak LR.}
    \label{tab:experiment_log}
\end{table}
